\section{Gabarito dos exercícios similares}

\subsection*{Parte 1}
\begin{description}[leftmargin=2.2em, itemsep=2pt, font=\bfseries]
\item[1.1] (a) Falso: $f(x)=x$ tem $f'=1>0$ mas $f<0$ para $x<0$. (b) Verdadeiro: $(f-g)'=0$, logo $f'=g'$. (c) Falso: $|x|$ é contínua em $0$ e não derivável em $0$.
\item[1.2] No instante $t=3$ min o volume está \emph{diminuindo} à taxa de $12$ L/min (sai mais água do que entra). Unidade: L/min.
\item[1.3] (a) $[x\sen x]_0^{\pi}=\pi\sen\pi-0=0$. \ (b) $0$ (derivada de uma constante). \ (c) $x\sen x$.
\item[1.4] Positiva: $\ln x>0$ em $(1,2]$, então a integral é a área sob uma curva positiva (propriedade 6).
\end{description}

\subsection*{Parte 2}
\begin{description}[leftmargin=2.2em, itemsep=2pt, font=\bfseries]
\item[2.1] $\dfrac{x^3-2x}{x^2}=x-\dfrac2x$ \;$\Rightarrow$\; $\dfrac{x^2}{2}-2\ln|x|+C$.
\item[2.2] $u^2-u-6$ \;$\Rightarrow$\; $\dfrac{u^3}{3}-\dfrac{u^2}{2}-6u+C$.
\item[2.3] $\sec^2\theta+1$ \;$\Rightarrow$\; $\tg\theta+\theta+C$.
\item[2.4] $f'(x)=2x+6x^2-4x^3+12$; \; $f(x)=-x^4+2x^3+x^2+12x+4$.
\item[2.5] $f(x)=x^2+x+C$; $f'(1)=3$; tangente $y=3(x-1)+2+C$ passa pela origem $\Rightarrow C=1$. \; $f(x)=x^2+x+1$.
\item[2.6] $v=-10\cos t+3\sen t+C_1$; $s=-10\sen t-3\cos t+C_1t+C_2$; $s(0)=0\Rightarrow C_2=3$; $s(2\pi)=2\pi C_1=12\Rightarrow C_1=6/\pi$. \; $s(t)=-10\sen t-3\cos t+\dfrac{6}{\pi}t+3$.
\end{description}

\subsection*{Parte 3}
\begin{description}[leftmargin=2.2em, itemsep=2pt, font=\bfseries]
\item[3.1] $\Delta x=\tfrac12$. $R_4=\tfrac12\big[\tfrac18+1+\tfrac{27}{8}+8\big]=6{,}25$ (super); $L_4=\tfrac12\big[0+\tfrac18+1+\tfrac{27}{8}\big]=2{,}25$ (sub). $x^3$ é crescente em $[0,2]$; exato $4$.
\item[3.2] $R_4=\tfrac14\big[\tfrac1{1{,}25}+\tfrac1{1{,}5}+\tfrac1{1{,}75}+\tfrac12\big]\approx0{,}6345$ (sub: $1/x$ é decrescente); $L_4\approx0{,}7595$ (super); exato $\ln2\approx0{,}6931$.
\item[3.3] $\Delta x=\tfrac2n$, $x_i=1+\tfrac{2i}{n}$, $f(x_i)=2+\tfrac{4i}{n}+\tfrac{4i^2}{n^2}$; $R_n=4+4\tfrac{n+1}{n}+\tfrac43\tfrac{(n+1)(2n+1)}{n^2}\to4+4+\tfrac83=\tfrac{32}{3}$. TFC: $[x^3/3+x]_1^3=12-\tfrac43=\tfrac{32}3$ \ok.
\item[3.4] (a) $\displaystyle\int_2^7(5x^3-4x)\dx$. \ (b) $\Delta x=\tfrac3n$, $x_i=\tfrac{3i}{n}$, $[0,3]$: $\displaystyle\int_0^3\sqrt{1+x}\dx=\Big[\tfrac23(1+x)^{3/2}\Big]_0^3=\tfrac23(8-1)=\tfrac{14}{3}$.
\item[3.5] $\Delta x=\tfrac2n$, $x_i=5+\tfrac{2i}{n}$: região sob $y=x^{10}$ de $x=5$ a $x=7$.
\item[3.6] (a) triângulos $\tfrac12\cdot1\cdot1+\tfrac12\cdot3\cdot3=5$. \ (b) semicírculo de raio $2$: $2\pi$. \ (c) reta zera em $x=2$: $-\tfrac12\cdot2\cdot4+\tfrac12\cdot1\cdot2=-4+1=-3$.
\item[3.7] $\displaystyle\lim_{n\to\infty}\sum_{i=1}^n e^{2i/n}\cdot\frac2n$.
\item[3.8] Em $[-1,1]$, $\sqrt{1+x^2}$ tem mínimo $1$ (em $x=0$) e máximo $\sqrt2$ (em $x=\pm1$); propriedade 8 com $b-a=2$.
\end{description}

\subsection*{Parte 4}
\begin{description}[leftmargin=2.2em, itemsep=2pt, font=\bfseries]
\item[4.1] (a) $\sqrt{x+x^3}$. \ (b) $\ln(e^x)\cdot e^x=xe^x$. \ (c) $\dfrac{\cos x}{1+\sen^2x}+\dfrac{\sen x}{1+\cos^2x}$.
\item[4.2] Derivando: $f(x)=3x^2$. Com $x=a$: $3=a^3\Rightarrow a=\sqrt[3]{3}$.
\item[4.3] $[x^2-x^4/4]_0^2=4-4=0$. A parte acima do eixo ($0\le x\le\sqrt2$, área $1$) cancela exatamente a parte abaixo ($\sqrt2\le x\le2$, área $1$).
\item[4.4] $\displaystyle\int_{-2}^0 2\dx+\int_0^2(4-x^2)\dx=4+\Big(8-\tfrac83\Big)=\tfrac{28}{3}$.
\item[4.5] $v=(t-3)(t+1)$: negativa em $[2,3)$, positiva em $(3,4]$. Deslocamento $=\tfrac23$ m; distância $=\tfrac53+\tfrac73=4$ m.
\item[4.6] Zero em $t=\tfrac53$. Deslocamento $=-\tfrac32$ m; distância $=\tfrac{25}{6}+\tfrac83=\tfrac{41}{6}\approx6{,}83$ m.
\item[4.7] $v(t)=\tfrac{t^2}{2}+4t+5>0$ sempre, então não há troca de sentido e distância $=$ deslocamento $=\big[\tfrac{t^3}{6}+2t^2+5t\big]_0^{10}=\tfrac{1250}{3}\approx416{,}7$ m.
\item[4.8] O volume de água no tanque em $t=4$ min ($25\,000$ L iniciais mais a variação líquida). Se $r<0$ em parte do intervalo, a integral continua sendo a variação \emph{líquida} (entradas menos saídas), e a expressão continua sendo o volume em $t=4$.
\item[4.9] $1/x^4$ não é contínua em $[-2,1]$ (explode em $x=0$); o TFC2 não se aplica e a integral não existe. O resultado negativo para um integrando positivo denuncia o erro.
\end{description}

\subsection*{Parte 5}
\begin{description}[leftmargin=2.2em, itemsep=2pt, font=\bfseries]
\item[5.1] $u=x^3$: $\tfrac13e^{x^3}+C$.
\item[5.2] $u=\ln x$: $\dfrac{(\ln x)^3}{3}+C$.
\item[5.3] $u=x^2+1$: $\tfrac13(x^2+1)^{3/2}+C$.
\item[5.4] $u=x^2$, limites $0\to1$: $\tfrac12[e^u]_0^1=\dfrac{e-1}{2}$.
\item[5.5] $u=\sen x$, limites $0\to1$: $[-\cos u]_0^1=1-\cos1\approx0{,}460$.
\item[5.6] $u=x^2$, limites $0\to\pi$: $\tfrac12[\sen u]_0^{\pi}=0$. A parte positiva de $x\cos(x^2)$ em $[0,\sqrt{\pi/2}]$ cancela a parte negativa em $[\sqrt{\pi/2},\sqrt\pi]$.
\item[5.7] $x^3$ é ímpar e $x^4\tg x$ é par$\times$ímpar $=$ ímpar; integral $=0$.
\item[5.8] $u=1+\tg x$, $du=\sec^2x\dx$: $\ln|1+\tg x|+C$.
\end{description}

\subsection*{Parte 6}
\begin{description}[leftmargin=2.2em, itemsep=2pt, font=\bfseries]
\item[6.1] Interseções $0$ e $2$; reta em cima: $\displaystyle\int_0^2(2x-x^2)\dx=4-\tfrac83=\tfrac43$.
\item[6.2] Interseções $-1,0,1$; por simetria $2\displaystyle\int_0^1(x-x^3)\dx=2\big(\tfrac12-\tfrac14\big)=\tfrac12$.
\item[6.3] $\displaystyle\int_0^2(2y-y^2)\,dy=4-\tfrac83=\tfrac43$.
\item[6.4] $\pi\displaystyle\int_1^2x^{-2}\dx=\pi\big[-\tfrac1x\big]_1^2=\tfrac{\pi}{2}$.
\item[6.5] $\displaystyle\int_0^12\pi x(x-x^2)\dx=2\pi\big(\tfrac13-\tfrac14\big)=\tfrac{\pi}{6}$.
\item[6.6] Cascas: $\displaystyle\int_0^4 2\pi x\sqrt x\dx=2\pi\cdot\tfrac25\cdot32=\tfrac{128\pi}{5}$. Anéis em $y$ ($0\le y\le2$, $R=4$, $r=y^2$): $\pi\displaystyle\int_0^2(16-y^4)\,dy=\pi\big(32-\tfrac{32}{5}\big)=\tfrac{128\pi}{5}$ \ok.
\end{description}
